\section{Conclusion}

In this work, we investigated a stochastic advection--diffusion equation driven by additive 
Gaussian noise. The model was formulated within a variational framework, allowing the 
application of standard tools from the theory of stochastic partial differential equations.

The mathematical analysis established the well-posedness of the problem. In particular, 
stability estimates were derived, and the existence and uniqueness of weak solutions were 
proved under suitable assumptions on the diffusion coefficient, velocity field, forcing 
term, and initial condition. Furthermore, a finite element approximation was introduced 
and its convergence properties were analyzed.

From the numerical point of view, several benchmark problems were considered. The finite 
element implementation was first validated on deterministic Poisson and advection--diffusion 
problems. Numerical experiments confirmed the expected second-order convergence in the 
\(L^2\)-norm with respect to the mesh size and first-order convergence with respect to the 
time step when the implicit Euler method was employed.

The stochastic model was then investigated through Monte Carlo simulations. The computed 
statistical quantities, including the mean energy and variance, demonstrated the stability 
of the numerical approximation and illustrated the propagation of uncertainty induced by 
the stochastic forcing. Additional Monte Carlo convergence studies showed the consistency 
of the statistical estimators as the number of realizations increased.

Overall, the numerical results are in excellent agreement with the theoretical analysis 
and confirm the reliability of the proposed finite element framework for stochastic 
advection--diffusion problems.

Future work will focus on several extensions of the present study. These include stochastic 
advection--diffusion--reaction equations, multiplicative noise models, adaptive finite 
element strategies, uncertainty quantification techniques, and optimal control problems 
governed by stochastic partial differential equations. Such developments are particularly 
relevant for environmental and climate-related applications involving uncertainty and 
random perturbations.
